% FAO-56 Kc curve for early-cycle (90-day) soybean planted on 1 December
% Source: Allen et al. (1998) -- FAO-56 Crop Evapotranspiration, Table 12
%         Steduto et al. (2012) -- FAO-66 Crop Yield Response to Water
% Effective rooting depth Z_r = 60 cm
% Available water capacity AWC = 120 mm
% Management allowed depletion MAD = 0.55 (FAO-56 default for soybean)
\begin{table}[!ht]
\centering
\caption{FAO-56 single-crop coefficient (\(K_c\)) for early-cycle soybean
(\textit{Glycine max} L.) planted on 1 December in MATOPIBA. Stage durations
and \(K_c\) values follow Allen \textit{et al.} (1998) Table 12 and Steduto
\textit{et al.} (2012) FAO-66.}
\label{tab:kc_soybean}
\begin{tabular}{lccc}
\toprule
\textbf{Phenological stage} & \textbf{Symbol} & \textbf{Duration (days)} & \textbf{\(K_c\)} \\
\midrule
Initial          & \(K_{c,\text{ini}}\) & 15 & 0.40 \\
Crop development & --                   & 15 & 0.80 \\
Mid-season       & \(K_{c,\text{mid}}\) & 40 & 1.15 \\
Late season      & --                   & 15 & 0.80 \\
Harvest          & \(K_{c,\text{end}}\) &  5 & 0.50 \\
\midrule
\textbf{Total cycle} & -- & \textbf{90} & -- \\
\bottomrule
\end{tabular}
\vspace{0.4em}\\
\footnotesize{Effective rooting depth \(Z_r = 60\) cm, available water
capacity AWC \(= 120\) mm, management allowed depletion MAD \(= 0.55\)
(FAO-56 Chapter 6 / FAO-66).}
\end{table}
